\section{partial derivatives}

It sometimes occurs that we have functions depending on multiple variables or
parameters. For example, the function $f(x,y) = x^2y+y^2$ is a function defined
on pairs of real numbers: $f\colon \RR^2\to \RR$. If we consider one of the two
variables, for instance the first variable $x$, as a fixed parameter, we can
identify the function $f$ with a function $g\colon \RR \to (\RR^\RR)$ that, for
each value of the parameter $x$, returns a function of the variable $y$:
\[
  g(x)(y) = f(x,y).
\]
Thus, for each fixed $x$, we can consider the function $h=g(x)$ which is a
function of the single variable $y$: $h(y) = g(x)(y) = f(x,y)$. We can then take
its derivative $h'(y)$ at any point $y$. The value obtained is called the
\emph{partial derivative}%
\mymargin{partial derivative}%
\index{derivative!partial}
of the function $f$ with respect to the variable $y$ evaluated at the point
$(x,y)$. To perform the derivative calculation, we can apply the usual
differentiation rules, taking care that if we are differentiating with respect
to $y$, the variable $x$ should be treated as if it were constant, because we
are effectively fixing $x$ before differentiating. If $f(x,y) = x^2y+y^2$, we
thus obtain%
\mynote{If we are not accustomed to thinking of $x$ as a fixed value, try substituting the variable $x$ with a different letter such as $c$ or $\pi$, which gives us the idea of a constant quantity.}:
$h'(y) = x^2 + 2y$.
If we now also consider $x$ as a variable, we obtain a new function of two
variables $(x,y)\mapsto (g(x))'(y) = h'(y)$ which can be written with the
following notations%
\mynote{%
Note that for functions of multiple variables, there is an inherent ambiguity in
the notations. Indeed, if I have an expression in two variables such as $x^2 y +
y^2$, it is clear that this represents a function of two variables, but it is
not entirely obvious which is the first variable and which is the second.
Normally, variables are indicated in alphabetical order starting from $x$... but
this is merely a convention. In principle, the expression $x^2y + y^2$ could
also represent a function of three variables $x,y,z$. Therefore, if I have an
expression that uses certain names for the variables, it is natural to use the
same names in the symbol used for the partial derivative, as in the notation
$D_y$. If, however, the variables do not have a name, it would be more sensible
to indicate the variable by its numerical position, as in the notation $D_2$.
There are cases where the notation can be very ambiguous, such as if I were to
write $\frac{\partial f(y,x)}{\partial y}$: am I differentiating $f$ with
respect to the first variable, or am I differentiating $f(x,y)$ with respect to
the second variable and then swapping the two variables?
}:
\begin{align*}
\frac{\partial f(x,y)}{\partial y} 
&= \frac{\partial f}{\partial y}(x,y)
= D_y f(x,y)
= f_y (x,y) \\
&= D_2 f(x,y) = f_{,2} (x,y) \\
&= h'(y)
\end{align*}

Similarly, we could take the partial derivative of $f$ with respect to the first
variable $x$. If $f(x,y) = x^2y+y^2$, we obtain
\[
\frac{\partial f(x,y)}{\partial x} = 2xy.
\]

Here, we are treating functions of multiple variables as if they were functions
of a single variable dependent on other parameters. This is because we are
conducting a course on analysis of functions of \emph{one} variable (real).
\index{one!variable}%
\index{function!of one variable}%
\index{variable!function of one}%
To treat all variables as a single multidimensional variable, several additional
considerations must be made, which are addressed in courses on analysis of
functions of \emph{multiple}
\index{multiple!variables}%
\index{function!of multiple variables}%
\index{variables!function of multiple}%
(real) variables. In such cases, concepts such as \emph{gradient} $\nabla f$,
\emph{derivative} $Df$, and \emph{differential} $df$ will be introduced, which
we will not address here.